\subsection{Historical context and the role of this reduction}
\label{sec:history}

Chatterjee, Thomas A.~Henzinger, and Jurdzi\'nski introduced
mean-payoff parity games to combine qualitative specifications with
quantitative performance requirements
\cite{ChatterjeeHenzingerJurdzinski2005}. They established the existence
of optimal strategies and showed that the player enforcing the
conjunction may need infinite memory. Bouyer, Markey, Olschewski,
and Ummels proved that the opposing player has an optimal positional
strategy and gave a simpler recursive value algorithm
\cite[Theorems~2 and~7]{BouyerMarkeyOlschewskiUmmels2011}.
Thus the strategy asymmetry and the recursive decomposition of these
games precede the present work.

Chatterjee and Doyen related mean-payoff parity games to energy parity
games, whose accumulated weight must remain above a fixed lower bound,
and established membership of the threshold problem in
$\mathrm{NP}\cap\mathrm{coNP}$~\cite{ChatterjeeDoyen2012}.
Chatterjee, Monika Henzinger, and Svozil improved the threshold
running time to $O(n^{d-1}mW)$, for $n$ vertices, $m$ edges,
$d\ge2$ priority levels, and $W$ the maximum of $1$ and the absolute weights
\cite[Theorem~6]{ChatterjeeHenzingerSvozil2017}.
This dependence is polynomial in the numerical weights but
exponential in the number of priorities. Daviaud, Jurdzi\'nski,
and Lazi\'c replaced the latter dependence by a quasipolynomial
factor, using strategy decompositions and progress measures combining
parity and energy information~\cite{DaviaudJurdzinskiLazic2018}.
Anand, Fijalkow, Goubault-Larrecq, Leroux, and Ohlmann subsequently
gave a separating-automaton construction for a related disjunction
of parity and mean-payoff objectives
with comparable pseudo-quasipolynomial complexity
\cite{AnandEtAl2021}. These bounds retain a numerical weight factor.

For parity games alone, Calude, Jain, Khoussainov, Li, and Stephan
obtained a quasipolynomial algorithm
\cite{CaludeEtAl2017}; Jurdzi\'nski and Lazi\'c developed the
succinct progress measures underlying the later mean-payoff parity
algorithm~\cite{JurdzinskiLazic2017}. The recursive ancestry relevant
here is Parys's modification of Zielonka's attractor recursion
\cite{Parys2019}, developed further by Lehtinen, Parys, Schewe,
and Wojtczak~\cite{LehtinenParysScheweWojtczak2022}.
Parys's Algorithm~2 and Lemma~6.1 already use separate size bounds
for the two players' dominions, two sequences of reduced-precision
calls surrounding one full-precision call, and the associated
dominion-preservation argument. Our algorithm adapts this structure
by repeatedly removing vertices where Max loses the ordinary
mean-payoff test. The strategy lemmas justify that adaptation for
the conjunction.

The strategy construction also has direct predecessors.
Increasing payoff phases interspersed with parity visits occur in
\cite{ChatterjeeHenzingerJurdzinski2005}; the even-priority combination
implication is stated in
\cite[Lemma~5]{BouyerMarkeyOlschewskiUmmels2011}.
The full version of that paper proves it with a phase of length $j$
in round $j$. The same scheduling appears in the threshold recursion
of~\cite[Appendix~C, full version]{ChatterjeeHenzingerSvozil2017}.
Daviaud, Jurdzi\'nski, and Lazi\'c also maintain a uniform bound on
finite-prefix costs in their strategy-decomposition proof, for their
strict mean-payoff convention
\cite[Sections~2.3--2.4]{DaviaudJurdzinskiLazic2018}.
Here we make the all-$\epsilon$ prefix guarantee explicit for a
single strategy at the non-strict threshold and use it to control
repeated restarts in the increasing-round construction against
history-dependent opponents.

The resulting reduction uses quasipolynomially many ordinary
mean-payoff queries of polynomial binary length. Substituting the
deterministic mean-payoff algorithm
\cite[Theorem~1.1 and Corollary~7.1]{OpenAI2026Deterministic}
therefore gives the full binary-length bound. This consequence uses
the present query bound together with that companion theorem; the
earlier recursive, strategy-combination, and dominion arguments supply
its methodological foundations.
